\documentclass[11pt,a4paper]{article}
\usepackage[margin=25mm]{geometry}
\usepackage{amsmath,amssymb,amsthm,booktabs}
\newtheorem{corollary}{Corollary}
\title{The order five and order eight entry}
\author{Human Compression project}\date{}
\begin{document}\maketitle
Let \(G=O'N\), of order \(460815505920\), and write \(x_H\) for the
number of ordered pairs \((r,s)\in H^2\) with
\(|r|=5,\ |s|=8,\ |rs|=2\). We use the acquired complete subgroup
roster and exact character evaluations, and the previously proved
result \emph{Unique overgroups from a normal subgroup of index two}.
In particular, that result proves that every subgroup in the eleven
classes
\[
 \mathcal R=\{7,8,30,33,34,58,59,105,106,133,186\}
\]
has exactly one actual maximal overgroup, in class 7. This assertion
concerns the subgroup classes themselves, so it applies to the present
pair orders without repeating the proof of the two normal-subgroup
intersections.

\begin{corollary}
There are \(470\) automorphism orbits of generating ordered pairs of
type \((5,8)\) in \(G\), with product of order two.
\end{corollary}
\begin{proof}
Let \(J=\langle r,s\rangle<G\). The complete upward closures of the
cyclic order-five class 574 and the two cyclic order-eight classes
559 and 560 have the following intersection:
\[
\begin{split}
1,7,8,11,12,15,16,18,19,23,24,30,32,33,34,58,59,60,61,\\
102,103,105,106,133,184,186.
\end{split}                                                    \tag{1}
\]
The character evaluations give \(x_{H_{11}}=x_{H_{12}}=0\).
Every subgroup of one of these two maximal groups also has zero
pairs of the requested type. In (1), rows \(11,18,19,32,60,61\)
occur below row 11, and rows \(12,184\) below row 12.
Removing them and \(G\) leaves precisely
\[
                 \mathcal R\ \cup\ \{15,16,23,24,102,103\}.
\]
The four rows \(15,16,23,24\) are maximal. Row 102 has only
row 23 as an immediate overgroup class, and row 103 only row 24.
There are \(58183776\) subgroups in each of classes 23 and 24,
and each contains eleven subgroups in the corresponding class
102 or 103. Each latter class has length \(640021536\). Since
\[
                   58183776\cdot11=640021536,
\]
incidence counting proves that the indicated parent is unique as
an actual subgroup. Together with the imported assertion for
\(\mathcal R\), this proves that every possible proper generated
subgroup has exactly one maximal overgroup.

It follows that summing the pair counts over all actual maximal
subgroups counts every nongenerating pair once. The only nonzero
maximal contributions are
\[
\begin{array}{c|r|r}
\text{row}&[G:H]&x_H\\\hline
7&2857239&2903040\\
15,16\text{ (each)}&30968784&119040\\
23,24\text{ (each)}&58183776&47520.
\end{array}
\]
The acquired ambient count is \(x_G=454364088837120\); the sum of
the maximal contributions is \(21197513272320\). Thus the number
of generating pairs is \(433166575564800\). An automorphism fixing
both generators is the identity, so the full automorphism group,
of order \(2|G|=921631011840\), acts freely. The quotient is
\[
                \frac{433166575564800}{921631011840}=470.
\]
There is no additional quotient by duality.
\end{proof}

This is the entry \((5,8)\) of Connor and Leemans,
\emph{Algorithmic enumeration of regular maps}, Table 1,
\emph{Ars Math. Contemp.} 10 (2016), 211--222.
The subgroup classification, character evaluations and automorphism
order remain explicitly borrowed premises. The earlier normal-subgroup
argument is reused; the additional work is the two immediate-parent
incidences and their application to this entry. No new finite group
enumeration, chirality conclusion, open-question solution, Lean
formalization or measured compression factor is claimed.
\end{document}
